\Legend
把答案看成一个阈值 $d$：要求任意距离小于 $d$ 的两点异色。

先考虑由两条边组成且三个点互不相同的路径，设最小长度为 $S$。当 $d>S$ 时，这条路径的三对点距离都小于 $d$（路径两端的最短距离至多为两边之和），不可能只用两种颜色让三点两两异色。因此答案不超过 $S$。计算 $S$ 时，对每个点找两条最短且通向\textbf{不同邻点}的边；平行边不可充当一条合法两边路径。

若 $d\le S$，任何长度小于 $d$ 的最短路径都不可能包含两条或更多边，否则前两条边就组成一条长度小于 $d$ 的合法两边路径。所以需要异色的点对恰好是边长小于 $d$ 的原图边。此时可行条件是这些边组成的图二分。

令 $B$ 为按边长从小到大加入边时，第一次出现奇环的边长。则 $d\le B$ 可以二分染色，$d>B$ 不行。答案因此为 $\min(S,B)$；若两者都不存在，则输出 $-1$。用带奇偶关系的并查集按边长扫描，若某条边的两端已被约束为同色，该边即给出 $B$。

主解时间 $O(m\log m+n)$，空间 $O(n+m)$。两边长度之和使用 64 位整数。参考解直接二分阈值并每次检查二分图，用于小规模交叉验证。
